\chapter{Thermodynamic potential of the Quark--Meson--Diquark model}
\label{app:qmd_potential}

This appendix collects the explicit expressions used for the
mean-field thermodynamic potential of the two-flavor
Quark--Meson--Diquark model. The expressions follow
Ref.~\cite{Andersen2024}. They are included here for reference in order
to keep the main text focused on the physical structure of the model.

The superconducting mean field is introduced by shifting the sigma field
and the third component of the diquark field according to
\begin{align}
\sigma &\rightarrow \phi_0+\tilde{\sigma}, \\
\Delta_3 &\rightarrow \Delta_0+\tilde{\Delta}_3 ,
\end{align}
where \(\phi_0\) and \(\Delta_0\) are real mean fields and the fluctuating
fields have vanishing expectation values. With the condensate chosen in
the third color direction, the tree-level thermodynamic potential is
\begin{align}
\Omega_0
&=
\frac{1}{2}m^2\phi_0^2
+
\left(m_\Delta^2-4\mu_3^2\right)\Delta_0^2
+
\frac{\lambda}{24}\phi_0^4
-
h\phi_0
+
\frac{\lambda_3}{12}\phi_0^2\Delta_0^2
+
\frac{\lambda_\Delta}{6}\Delta_0^4 ,
\label{eq:app_qmd_tree_potential}
\end{align}
where
\begin{equation}
\mu_3=\frac{1}{2}\left(\mu_{ur}+\mu_{dg}\right).
\end{equation}

For the explicit evaluation of the quark determinant, we use the
Nambu--Gorkov basis defined in Eq.~\eqref{eq:qmd_nambu_gorkov_spinor}.
The corresponding inverse propagator is
\begin{equation}
D^{-1}(p)
=
\begin{pmatrix}
\slashed{p}-g\phi_0+\gamma^0\hat{\mu}
&
i g_\Delta \tau_2 \lambda_2 \gamma^5 \Delta_0
\\
i g_\Delta \tau_2 \lambda_2 \gamma^5 \Delta_0
&
\slashed{p}-g\phi_0-\gamma^0\hat{\mu}
\end{pmatrix}.
\label{eq:app_qmd_inverse_propagator}
\end{equation}
Here the tensor \(\epsilon_{3ab}\) has been replaced by the Gell-Mann
matrix \(i\lambda_2\), corresponding to the chosen color orientation of the
condensate.

The quasiparticle energies entering the one-loop potential are
\begin{align}
E_{ub}^{\pm}(p) &= E(p)\pm\mu_{ub}, \\
E_{db}^{\pm}(p) &= E(p)\pm\mu_{db}, \\
E_{\Delta^\pm}^{\pm}(p) &= E_\Delta^\pm(p)\pm\delta\mu,
\end{align}
where
\begin{align}
E(p) &= \sqrt{p^2+g^2\phi_0^2}, \\
E_\Delta^\pm(p)
&=
\sqrt{\left(E(p)\pm\bar{\mu}\right)^2
+
g_\Delta^2\Delta_0^2},
\\
\bar{\mu}
&=
\frac{1}{2}\left(\mu_{ur}+\mu_{dg}\right)
=
\frac{1}{2}\left(\mu_{ug}+\mu_{dr}\right),
\\
\delta\mu
&=
\frac{1}{2}\left(\mu_{dg}-\mu_{ur}\right)
=
\frac{1}{2}\left(\mu_{dr}-\mu_{ug}\right).
\end{align}
For the common quark chemical potential used in Ref.~\cite{Andersen2024},
\(\delta\mu=0\), \(\mu_3=\bar{\mu}\), and \(\mu_B=3\bar{\mu}\).

The full renormalized one-loop thermodynamic potential can be written as
\begin{align}
\Omega_{0+1}
&=
\frac{3}{4}m_\pi^2 f_\pi^2
\left[
1-
\frac{12m_q^2}{(4\pi)^2f_\pi^2}
m_\pi^2F'(m_\pi^2)
\right]
\frac{\phi_0^2}{f_\pi^2}
\nonumber \\
&\quad
+
\frac{2m_q^4}{(4\pi)^2}
\left[
\frac{9}{2}
+
\log\!\left(\frac{m_q^2}{g_0^2\phi_0^2}\right)
+
2\log\!\left(
\frac{m_q^2}{g_0^2\phi_0^2+g_{\Delta,0}^2\Delta_0^2}
\right)
\right]
\frac{\phi_0^4}{f_\pi^4}
\nonumber \\
&\quad
-
\frac{1}{4}m_\sigma^2 f_\pi^2
\left\{
1+
\frac{12m_q^2}{(4\pi)^2f_\pi^2}
\left[
\left(1-\frac{4m_q^2}{m_\sigma^2}\right)F(m_\sigma^2)
+
\frac{4m_q^2}{m_\sigma^2}
-
F(m_\pi^2)
-
m_\pi^2F'(m_\pi^2)
\right]
\right\}
\frac{\phi_0^2}{f_\pi^2}
\nonumber \\
&\quad
+
\frac{1}{8}m_\sigma^2 f_\pi^2
\left\{
1+
\frac{12m_q^2}{(4\pi)^2f_\pi^2}
\left[
\left(1-\frac{4m_q^2}{m_\sigma^2}\right)F(m_\sigma^2)
-
F(m_\pi^2)
-
m_\pi^2F'(m_\pi^2)
\right]
\right\}
\frac{\phi_0^4}{f_\pi^4}
\nonumber \\
&\quad
-
\frac{1}{8}m_\pi^2 f_\pi^2
\left[
1-
\frac{12m_q^2}{(4\pi)^2f_\pi^2}
m_\pi^2F'(m_\pi^2)
\right]
\frac{\phi_0^4}{f_\pi^4}
\nonumber \\
&\quad
-
m_\pi^2 f_\pi^2
\left[
1-
\frac{12m_q^2}{(4\pi)^2f_\pi^2}
m_\pi^2F'(m_\pi^2)
\right]
\frac{\phi_0}{f_\pi}
\nonumber \\
&\quad
+
\left\{
m_{\Delta,0}^2
-
4\bar{\mu}^2
\left[
1+
\frac{4g_{\Delta,0}^2}{(4\pi)^2}
\left(
\log\!\left(
\frac{m_q^2}
{g_0^2\phi_0^2+g_{\Delta,0}^2\Delta_0^2}
\right)
-
F(m_\pi^2)
-
m_\pi^2F'(m_\pi^2)
\right)
\right]
\right\}
\Delta_0^2
\nonumber \\
&\quad
+
\frac{\lambda_{3,0}}{12}\phi_0^2\Delta_0^2
+
\frac{\lambda_{\Delta,0}}{6}\Delta_0^4
+
\frac{12g_0^2g_{\Delta,0}^2}{(4\pi)^2}\phi_0^2\Delta_0^2
+
\frac{6g_{\Delta,0}^4}{(4\pi)^2}\Delta_0^4
\nonumber \\
&\quad
+
\frac{4g_{\Delta,0}^4}{(4\pi)^2}
\left[
\log\!\left(
\frac{m_q^2}
{g_0^2\phi_0^2+g_{\Delta,0}^2\Delta_0^2}
\right)
-
F(m_\pi^2)
-
m_\pi^2F'(m_\pi^2)
\right]\Delta_0^4
\nonumber \\
&\quad
+
\frac{4g_0^2g_{\Delta,0}^2}{(4\pi)^2}
\left[
\log\!\left(
\frac{m_q^2}
{g_0^2\phi_0^2+g_{\Delta,0}^2\Delta_0^2}
\right)
-
F(m_\pi^2)
-
m_\pi^2F'(m_\pi^2)
\right]\phi_0^2\Delta_0^2
\nonumber \\
&\quad
+
\Omega_{1,\mathrm{num}}
+
\Omega^{\bar{\mu},T}.
\label{eq:app_qmd_full_potential}
\end{align}

The function \(F(p^2)\) and its derivative \(F'(p^2)\) are defined in
Appendix~\ref{app:loop_integrals}. In the present appendix they are
evaluated with the vacuum constituent quark mass \(m_q=g_0 f_\pi\).
\begin{equation}
\Omega^{\bar{\mu},T}
=
-4T
\int_p
\left\{
\log\!\left[1+e^{-\beta(E\pm\bar{\mu})}\right]
+
2\log\!\left[1+e^{-\beta E_\Delta^\pm}\right]
\right\}.
\label{eq:app_qmd_temperature_term}
\end{equation}
In the zero-temperature calculations considered in the main text this term
is evaluated in the limit \(T\rightarrow0\).

The finite residual contribution \(\Omega_{1,\mathrm{num}}\) is
\begin{align}
\Omega_{1,\mathrm{num}}
&=
-4
\int_p
\bigg[
E_\Delta^+(p)
+
E_\Delta^-(p)
-
2\sqrt{
p^2+g^2\phi_0^2+g_\Delta^2\Delta_0^2
}
\nonumber \\
&\hspace{3.0cm}
-
\frac{
\bar{\mu}^2 g_\Delta^2\Delta_0^2
}{
\left(
p^2+g^2\phi_0^2+g_\Delta^2\Delta_0^2
\right)^{3/2}
}
\bigg].
\label{eq:app_qmd_omega_num}
\end{align}
This integral is finite. It arises after subtracting the ultraviolet
divergent part of the one-loop quark determinant. Unlike the other terms
in Eq.~\eqref{eq:app_qmd_full_potential}, it is not available as a simple
closed analytic expression and must be treated separately in the numerical
implementation.
